\section{Bounded overlap does not imply coverage}
\label{sec:counterexamples}

Bounded overlap provides an upper-norm inequality.  It does not
provide a lower frame bound, surjectivity, or coefficient feasibility.

\begin{lemma}[What overlap actually controls]
\label{lem:overlap-upper-bound}
Put \(S_j=\{\omega\in\Omega:I_{\omega j}=1\}\).  For the weighted
incidence matrix \eqref{eq:weighted-incidence},
\begin{equation}
 \|Bc\|_2^2
 \le
 \Delta(I)
 \sum_{\omega\in\Omega}\sum_{j=1}^m
 I_{\omega j}|g_{\omega j}c_j|^2.
 \label{eq:overlap-upper-bound}
\end{equation}
In particular, if \(|g_{\omega j}|\le G\), then
\[
 \|Bc\|_2^2
 \le
 \Delta(I)G^2
 \sum_{j=1}^m |c_j|^2|S_j|.
\]
\end{lemma}

\begin{proof}
For each row \(\omega\), Cauchy--Schwarz gives
\[
 \left|
 \sum_j I_{\omega j}g_{\omega j}c_j
 \right|^2
 \le
 \left(\sum_jI_{\omega j}\right)
 \sum_jI_{\omega j}|g_{\omega j}c_j|^2.
\]
Sum over \(\omega\).  The last displayed specialization follows by
interchanging the sums and using
\(\sum_\omega I_{\omega j}=|S_j|\).
\end{proof}

\begin{remark}
The right side of \eqref{eq:overlap-upper-bound} is an upper synthesis
bound.  No reverse inequality follows from \(\Delta(I)<\infty\):
columns may be linearly dependent, miss directions, or conflict with
the target weights.
\end{remark}

\subsection{A disjoint full-support family with fixed residual}

\begin{theorem}[Overlap one, but no asymptotic cover]
\label{thm:partition-counterexample}
For every \(n\ge1\), there is an unweighted incidence matrix
\(B_n\in\{0,1\}^{2n\times n}\) and a positive target
\(w_n\in\R^{2n}\) such that:
\begin{enumerate}[label=(\roman*),leftmargin=2.2em]
\item every atom lies in exactly one block, so
  \(\Delta(I)=1\) and the blocks partition the full support;
\item no signed, complex, or nonnegative exact cover exists;
\item the exact least-squares defect is
  \[
  \dist(w_n,\ran B_n)=\sqrt{\frac n2};
  \]
\item the relative defect is the fixed constant
  \[
  \frac{\dist(w_n,\ran B_n)}{\|w_n\|_2}
  =\frac1{\sqrt{10}}.
  \]
\end{enumerate}
\end{theorem}

\begin{proof}
Let
\[
 b_j=e_{2j-1}+e_{2j},
 \qquad
 w_n=\sum_{j=1}^n(e_{2j-1}+2e_{2j}).
\]
The columns have disjoint supports and every row is present once.
For any \(c=(c_j)\),
\[
 \|w_n-B_nc\|_2^2
 =
 \sum_{j=1}^n
 \left(|1-c_j|^2+|2-c_j|^2\right).
\]
Each summand is uniquely minimized at \(c_j=3/2\), where it equals
\(1/2\).  The minimizer is nonnegative, so imposing
\(c\ge0\) does not improve or worsen the minimum.  Hence the squared
defect is \(n/2\), whereas
\(\|w_n\|_2^2=5n\).  The positive defect excludes every exact cover.
\end{proof}

\begin{corollary}[No approximate-cover consequence]
\label{cor:no-approx-from-overlap}
Even the conjunction of bounded overlap, full set coverage, positive
target weights, and a disjoint block partition does not imply
\[
 \dist(w_X,\ran B_X)=o(\|w_X\|_2).
\]
\end{corollary}

\subsection{Signed feasibility does not imply allowed feasibility}

\begin{theorem}[A sharp two-column cone obstruction]
\label{thm:cone-counterexample}
Let
\[
 B=
 \begin{pmatrix}
 1&0\\
 1&1
 \end{pmatrix},
 \qquad
 w=
 \begin{pmatrix}
 1\\0
 \end{pmatrix}.
\]
Then:
\begin{enumerate}[label=(\roman*),leftmargin=2.2em]
\item every row occurs in a block and the overlap degree is two;
\item \(Bc=w\) has the unique signed solution \(c=(1,-1)^{\mathsf T}\);
\item no nonnegative cover exists;
\item
  \[
  \dist(w,B\R_+^2)=\frac1{\sqrt2}.
  \]
\end{enumerate}
The vector \(y=(-1,1)^{\mathsf T}\) is a Farkas certificate:
\[
 B^{\mathsf T}y=(0,1)^{\mathsf T}\ge0,
 \qquad
 \langle y,w\rangle=-1<0.
\]
\end{theorem}

\begin{proof}
The signed equation gives \(c_1=1\) from the first row and
\(c_1+c_2=0\) from the second.  For \(c_1,c_2\ge0\),
\[
 \|w-Bc\|_2^2
 =
 (1-c_1)^2+(c_1+c_2)^2.
\]
At fixed \(c_1\), the minimum in \(c_2\ge0\) occurs at \(c_2=0\).
The remaining quadratic is minimized at \(c_1=1/2\), with value
\(1/2\).  The displayed Farkas signs are immediate.
\end{proof}

\subsection{Native-key collapse can create a false cover}

Let \(A:\K^\Omega\to\K^\Gamma\) aggregate literal atoms into coarse
classes, for example by summing all atoms with equal numerical
determinant.

\begin{proposition}[Coarse feasibility need not lift]
\label{prop:coarse-does-not-lift}
The coarse equation
\begin{equation}
 ABc=Aw
 \label{eq:coarse-cover}
\end{equation}
is necessary but not sufficient for the literal equation \(Bc=w\).
The failure can occur with one block, one coarse key, and overlap one.
\end{proposition}

\begin{proof}
Take two literal atoms with the same coarse key,
\[
 A=\begin{pmatrix}1&1\end{pmatrix},
 \qquad
 B=e_1,
 \qquad
 w=e_2.
\]
Then \(AB=Aw=1\), so \(c=1\) solves the coarse equation.  But
\(Bc=e_1\ne e_2=w\), and the literal defect is
\(\|e_2-e_1\|_2=\sqrt2\).  The two atoms may have the same determinant
and different provenance; aggregation has hidden exactly that
difference.
\end{proof}

\begin{remark}[Union-frame rule]
If a candidate column emits atoms outside \(\supp w\), the correct
universe is the union of target and column supports, with target
coefficient zero on the extra rows.  Deleting those rows can turn
leakage into a false exact cover, just as
\cref{prop:coarse-does-not-lift} turns a literal mismatch into a
coarse identity.
\end{remark}

\begin{example}[Overlap two and inconsistent multiplicity]
\label{ex:three-row-overlap}
Let
\[
 B=
 \begin{pmatrix}
 1&0\\
 1&1\\
 0&1
 \end{pmatrix},
 \qquad
 w=(1,1,1)^{\mathsf T}.
\]
The first and third rows force \(c_1=c_2=1\), while the middle row
would then equal \(2\).  Thus there is no exact cover although every
atom is covered and the overlap is two.  The least-squares solution is
\(c_1=c_2=2/3\), with residual
\((1,-1,1)^{\mathsf T}/3\) and relative defect \(1/3\).
It gives a three-row binary example of why an ad hoc division by
overlap multiplicity cannot replace the literal system.
\end{example}
